\section{Conclusion}
In this work we presented an open-source, catalog-agnostic Python package that crossmatches a stellar catalog against a planetary catalog and enriches the matched planets with stellar and planetary parameters. Applied to HPICxLIFEcat, the pipeline was run against both the NASA Exoplanet Archive and Exo-MerCat, prioritizing identifier-based matching, using alternate identifiers from SIMBAD and Exo-MerCat, over a proper-motion-aware coordinate-based fallback.

Identifier-based matching proved to be the dominant matching method by number of matched planets, with coordinate-only matches remaining the edge case. Calibrating the coordinate thresholds on ID-matched stars showed that the proper-motion-aware search radius with a base tolerance of $10''$ recovers all matches against NEA, and that the flat $50''$ fallback radius is sufficient when proper motion is unavailable; in both cases the sparse on-sky density of exoplanet hosts leads to no predicted false positives. The Exo-MerCat matches form a near superset of the NEA matches.

After enrichment from a priority-ordered set of catalog sources and physically motivated fallback derivations, the combined catalog yields 20 rocky, temperate planet candidates around HPICxLIFEcat stars (7 confirmed and 13 uncertain), out of 83 across the entire Exo-MerCat catalog. These planets are hosted almost exclusively by M dwarfs, reflecting the abundance of cool, low-mass stars in the immediate solar neighborhood that dominates the LIFE target list. This list provides a set of already-known targets relevant to LIFE's search phase, and, owing to the catalog-agnostic design of the package, can be regenerated quickly as the stellar target-selection catalog and the planetary catalogs evolve.
